\documentclass[12pt,a4paper]{article}
% \usepackage[utf8]{inputenc} % fontspecと併用時は不要
\usepackage[japanese]{babel}
\usepackage{amsmath}
\usepackage{amsfonts}
\usepackage{amssymb}
\usepackage{geometry}
\usepackage{graphicx}
\usepackage{enumitem}
\usepackage{fancyhdr}
\usepackage{verbatim}

% 日本語改行設定
\usepackage{xeCJK}
\usepackage{indentfirst}
\setlength{\parindent}{1em}

% 日本語フォント設定
\usepackage{fontspec}
\setmainfont{Hiragino Mincho Pro}
\setsansfont{Hiragino Sans}
\setmonofont{Menlo}
\setCJKmainfont{Hiragino Mincho Pro}
\setCJKsansfont{Hiragino Sans}
\setCJKmonofont{Menlo}
% ページ設定
\geometry{margin=1.5cm}
\setlength{\headheight}{14.70604pt}
\pagestyle{fancy}
\fancyhf{}
\fancyhead[L]{プログラミング応用 第五回}
\fancyhead[R]{演習問題}
\fancyfoot[R]{\ifnum\value{page}=1\small（裏面に続く）\fi}
\fancyfoot[C]{\thepage}
\renewcommand{\headrulewidth}{0.4pt}
\renewcommand{\footrulewidth}{0.4pt}

% 問題番号のカウンター
\newcounter{problem}
\setcounter{problem}{0}

% シンプルな問題環境
\newenvironment{problem}[1][]{
    \refstepcounter{problem}
    \vspace{1em}
    \noindent
    \textbf{問\arabic{problem}} #1
    %\vspace{0.5em}
    %\par
}{
    \vspace{1em}
}

\begin{document}

% ヘッダー情報
\begin{center}
    {\Large \textbf{プログラミング応用 第五回}}\\[0.5em]
    {\large \textbf{演習問題}}\\[1em]
\end{center}

\noindent\textbf{注意事項}

\begin{itemize}[leftmargin=1em]
    \item 解答はpdf形式でLMSにアップロードしてください (例えば紙に手書きしてスキャンしたものや \TeX をコンパイルしたものなど)
    \item 締切：2025年11月\textbf{18日（火）}日本時間13時
    \item 気をつけていますが\textbf{問題文に不備がある可能性はあります}ので, 気軽に問い合わせてください (SlackのDMも常時受付ます).
    \item 中々解けない問題についてはヒントも出しますので, 気軽に質問してください.
    \item なお、アルゴリズムの記述はこれまでの講義資料を参考にしてください. \textbf{帰着に関する問題は多項式時間であることが明確に分かれば良いので, 計算モデルに沿って厳密に議論しなくてよい (帰着の本質的なアイデアを評価する).}
\end{itemize}

\begin{problem}
    以下の各問いに答えよ.
    \begin{enumerate}[label=(\arabic*)]
        \item ハミルトンパス問題をハミルトン閉路問題にカープ帰着せよ.
        \item 次の判定問題を$\Pi_{\mathrm{path}}$とする: 重みつき無向グラフ $G=(V,E,w)$ (ただし$w\colon E \to \mathbb{Z}$ は負の重みも許される), 二つの頂点 $s,t\in V$, そして整数 $c\in \mathbb{Z}$ に対して, $s$から$t$への長さ$c$以下のパスが存在するならばYes, そうでないならばNoを出力せよ. ここで, パスとは路であって\textbf{同じ頂点を二度以上通らない}ものをいい, パスの長さとはそのパスに含まれる辺の重みの総和である. ハミルトン閉路問題を$\Pi_{\mathrm{path}}$にカープ帰着せよ.
    \end{enumerate}
\end{problem}

\begin{problem}
    ナップザック問題に関する以下の二つの判定問題を考える:
    \begin{itemize}    
        \item \textbf{ナップザック判定問題}: $n$個のアイテムがあり, 各アイテム$i$は価値 $v_i\in\mathbb{Z}_{\ge 0}$ と重み $w_i\in\mathbb{Z}_{\ge 0}$ を持つ. 各$(v_i,w_i)$ ($i=1,\dots,n$), ナップザックの容量 $W\in\mathbb{Z}_{\ge 0}$, そして閾値 $K\in\mathbb{Z}_{\ge 0}$ が与えられる. 容量を超えない範囲で選択したアイテムの価値の総和が$K$以上となる組み合わせが存在するならばYes, そうでなければNoを出力せよ. ただし, \textbf{各アイテムはいくつでも詰め込むことができる}.
        \item \textbf{01ナップザック判定問題}: 各アイテムは\textbf{高々一つまで}詰め込めるとしたときのナップザック判定問題. すなわち, $n$個のアイテムがあり, 各アイテム$i$は価値 $v_i\in\mathbb{Z}_{\ge 0}$ と重み $w_i\in\mathbb{Z}_{\ge 0}$ を持つ. 各$(v_i,w_i)$ ($i=1,\dots,n$), ナップザックの容量 $W\in\mathbb{Z}_{\ge 0}$, そして閾値 $K\in\mathbb{Z}_{\ge 0}$ が与えられる. 容量を超えない範囲で選択したアイテムの価値の総和が$K$以上となる組み合わせが存在するならばYes, そうでなければNoを出力せよ. ただし, \textbf{各アイテムは高々一つまで詰め込むことができる}.
    \end{itemize}

    ナップザック判定問題を01ナップザック判定問題にカープ帰着せよ. 入力で与えられる整数の桁数に注意すること.
\end{problem}

\begin{problem}
$n$人でゲームを行う ($n\geq 2$). このゲームは二つのチームに分かれて行うゲームである.
各プレイヤー$i$には戦力$w_i \in \mathbb{Z}_{\ge 0}$ が割り当てられており, チームの戦力はそのチームに属するプレイヤーの戦力の総和と定義する.
$n$人のプレイヤーを二つのチームにチーム分けすることを考える.
全てのプレイヤーはどちらか一方のみのチームに属するとする (即ち, 両方のチームに所属していたり, どちらのチームにも属さないプレイヤーはいないとする).
なお, 両チームの人数が同じである必要はない.
上手くチーム分けすることによって, 両チームの戦力が同じ値になるようにできるならばYes, そうでなければNoを出力する判定問題を\textbf{戦力均衡分配問題}と呼ぶ.

戦力均衡分配問題を01ナップザック判定問題にカープ帰着せよ.
\end{problem}

\newpage

\begin{problem}
    計算量クラスに関する以下の問題に答えよ.
    \begin{enumerate}[label=(\arabic*)]
        \item $\mathtt{P}\subseteq\mathtt{NP}$ を示せ.
        \item $\Pi_1\le_{\mathrm{Karp}} \Pi_2$ かつ $\Pi_2\in\mathtt{NP}$ ならば $\Pi_1\in\mathtt{NP}$ を示せ.
    \end{enumerate}
\end{problem}

\begin{problem}
以下の二つは同値であることを示せ.
\begin{itemize}
    \item あるNP完全な判定問題 $\Pi$ に対して, $\Pi \in \mathtt{P}$
    \item $\mathtt{P}=\mathtt{NP}$
\end{itemize}
\end{problem}

\end{document}
